% !TeX root = ../navier-stokes-manuscript.tex
% ============================================================
% 2.1 Vortex Stretching
% ============================================================

Axial stretching lengthens one narrow material vortex-tube segment while its material volume stays fixed, so its cross-section contracts, its rotating fluid moves toward the axis, and the interior spins faster.
The velocity field carrying this segment obeys the momentum balance
\[
  \underbrace{\partial_t u}_{\text{local acceleration}}
  +
  \underbrace{(u\cdot\nabla)u}_{\text{nonlinear transport}}
  =
  \underbrace{-\nabla p}_{\text{pressure}}
  +
  \underbrace{\nu\Delta u}_{\text{viscous diffusion}}.
\]
Taking its curl removes pressure and turns nonlinear transport into material vorticity advection and stretching, where viscosity acts at the same time as the segment narrows.

\subsection{Vortex Stretching}\label{sub:ThePerfectlyStretchingVortex}

The segment is locally axisymmetric.
Its strain is locally uniform across the material segment, and that strain is aligned with the segment's axis.
If \(e\) denotes this axis and \(L(t)\) the segment's length, the aligned extension has the form
\[
  S:=\frac12\bigl(\nabla u+(\nabla u)^{\mathsf T}\bigr),
  \qquad
  Se=\sigma(t)e,
  \qquad
  \sigma(t)>0,
  \qquad
  \frac{\dot L(t)}{L(t)}
  =
  e\cdot Se
  =
  \sigma(t).
\]
As \(L(t)\) grows, incompressibility keeps the material volume \(\mathcal V\) fixed.
With effective transverse area \(A(t):=\mathcal V/L(t)\) and effective radius \(r_{\mathrm{eff}}(t):=\sqrt{A(t)/\pi}\), the same extension gives
\[
  \nabla\cdot u=0,
  \qquad
  \frac{d}{dt}(A L)=0,
  \qquad
  \frac{\dot A}{A}=-\sigma(t),
  \qquad
  \frac{\dot r_{\mathrm{eff}}}{r_{\mathrm{eff}}}
  =
  -\frac{\sigma(t)}{2}.
\]
The contracting radius carries rotating fluid inward.
Writing \(\omega:=\nabla\times u\), its vorticity evolves by
\[
  \frac{D\omega}{Dt}
  =
  S\omega+\nu\Delta\omega,
  \qquad
  \frac{D}{Dt}
  :=
  \partial_t+u\cdot\nabla,
\]
and under the alignment \(\omega=|\omega|e\), the stretching and viscous actions remain together:
\[
  S\omega
  =
  \sigma(t)\omega,
  \qquad
  \frac12\frac{D|\omega|^2}{Dt}
  =
  \sigma(t)|\omega|^2
  +
  \nu\,\omega\cdot\Delta\omega.
\]
The inward-moving fluid spins faster while viscosity acts simultaneously, and squared vorticity grows only on intervals where this complete right-hand side is positive.
The identity proved in \Cref{prop:ClassicalKineticEnergyIdentity} controls the velocity in \(L^2\), while vorticity lies in the antisymmetric part of its gradient:
\[
  \norm{u(t)}_2
  \leq
  \norm{u_0}_2
  <\infty,
  \qquad
  \left|\frac12\bigl(\nabla u-(\nabla u)^{\mathsf T}\bigr)\right|_{\mathrm F}
  =
  \frac{|\omega|}{\sqrt2}
  \leq
  |\nabla u|_{\mathrm F}.
\]
Finite kinetic energy does not prevent that rotational gradient from concentrating inside the narrowing segment.

\divider{\NavierStokes{} Scaling}

At finer width, the vorticity remains coupled to velocity, pressure, transport, and viscosity, so the full field changes scale together.
This unresolved concentration motivates the comparison without making the energy bound a proof of the scaling symmetry.
At a time \(t_0<T_*\), set \(\ell:=r_{\mathrm{eff}}(t_0)\).
For \(\lambda>1\), rescale the velocity and pressure together by
\[
  u_\lambda(x,t)
  :=
  \lambda u(\lambda x,\lambda^2t),
  \qquad
  p_\lambda(x,t)
  :=
  \lambda^2p(\lambda x,\lambda^2t).
\]
On the corresponding interval, \(u_\lambda\) is another \NavierStokes{} solution, and at time \(\lambda^{-2}t_0\) its corresponding segment has effective radius \(\lambda^{-1}\ell\); it is not a later material state of the opening segment.
The transformation of each momentum term is
\[
\begin{aligned}
  \partial_tu_\lambda(x,t)
  &=
  \lambda^3(\partial_tu)(\lambda x,\lambda^2t),
  \\
  (u_\lambda\cdot\nabla)u_\lambda(x,t)
  &=
  \lambda^3\bigl((u\cdot\nabla)u\bigr)(\lambda x,\lambda^2t),
  \\
  \nabla p_\lambda(x,t)
  &=
  \lambda^3(\nabla p)(\lambda x,\lambda^2t),
  \\
  \nu\Delta u_\lambda(x,t)
  &=
  \lambda^3\nu(\Delta u)(\lambda x,\lambda^2t).
\end{aligned}
\]
Every term gains the same cubic factor, so the equation stays scale matched even though its Sobolev norms change with order.

\divider{Critical Height}

Let \(\Lambda:=(-\Delta)^{1/2}\).
The Fourier transform of the rescaled velocity is
\[
  \widehat{u_\lambda}(\xi,t)
  =
  \lambda^{-2}\widehat u(\lambda^{-1}\xi,\lambda^2t),
\]
and a change of variables gives the order-dependent law
\[
  \norm{u_\lambda(t)}_{\dot H^s}^2
  =
  \lambda^{2s-1}
  \norm{u(\lambda^2t)}_{\dot H^s}^2.
\]
At \(s=0\), kinetic energy loses one power of \(\lambda\); at \(s=1\), the first derivatives gain one.
At \(s=\tfrac12\), the exponent vanishes, and the corresponding half-derivative quantity
\[
  H(t)
  :=
  \frac12\norm{u(t)}_{\dot H^{1/2}}^2
  =
  \frac12\norm{\Lambda^{1/2}u(t)}_2^2.
\]
is the critical height.
Let \(H_\lambda\) denote that same quantity evaluated on \(u_\lambda\); then
\[
  \norm{u_\lambda(t)}_2^2
  =
  \lambda^{-1}\norm{u(\lambda^2t)}_2^2,
  \qquad
  H_\lambda(t)=H(\lambda^2t),
  \qquad
  \norm{\nabla u_\lambda(t)}_2^2
  =
  \lambda\norm{\nabla u(\lambda^2t)}_2^2.
\]
The middle quantity is neutral under scale, but that neutrality leaves the sign of its change in time unresolved.

\divider{Nonlinear Production versus Viscous Dissipation}

Nonlinear transport changes \(H\) through the signed production
\[
  \mathcal P_{1/2}[u](t)
  :=
  -\operatorname{Re}
  \pair
    {\Lambda^{1/2}u(t)}
    {\Lambda^{1/2}\bigl((u(t)\cdot\nabla)u(t)\bigr)}_{L^2}.
\]
Incompressibility removes the pressure pairing, while viscosity subtracts the nonnegative dissipation \(\nu\norm{\Lambda^{3/2}u}_2^2\).
The exact balance is
\[
  H'(t)
  +
  \nu\norm{\Lambda^{3/2}u(t)}_2^2
  =
  \mathcal P_{1/2}[u](t).
\]
Its signed right-hand side decides whether \(H\) rises or falls after dissipation is accounted for.
Under the same rescaling, production and dissipation satisfy
\[
  \mathcal P_{1/2}[u_\lambda](t)
  =
  \lambda^2\mathcal P_{1/2}[u](\lambda^2t),
  \qquad
  \nu\norm{\Lambda^{3/2}u_\lambda(t)}_2^2
  =
  \lambda^2\nu\norm{\Lambda^{3/2}u(\lambda^2t)}_2^2.
\]
Their common quadratic scaling preserves the balance but supplies no time in which the nonlinear field forms the next smaller structure.

\divider{The Heat Clock}

Viscosity has a separate linear comparison clock.
For the heat equation \(v_t=\nu\Delta v\),
\[
  \widehat v(\xi,t+\delta t)
  =
  e^{-\nu|\xi|^2\delta t}\widehat v(\xi,t).
\]
A scale-localized structure of width \(r\) occupies Fourier frequencies \(|\xi|\simeq r^{-1}\), so its linear viscous change becomes order one when
\[
  \nu|\xi|^2\delta t\simeq1,
\]
and the resulting comparison time is
\[
  \tau_\nu(r)
  :=
  \frac{r^2}{\nu}.
\]
This linear clock does not establish the lifetime of a nonlinear structure.
Its quadratic contraction is nonetheless exact:
\[
  \tau_\nu(\lambda^{-1}r)
  =
  \lambda^{-2}\tau_\nu(r).
\]
Dyadic widths thus produce geometrically shrinking comparison times.

\Needspace{18\baselineskip}
\divider{The Zeno Cascade}

For
\[
  r_n:=2^{-n}r_0,
  \qquad
  \tau_n:=\frac{r_n^2}{\nu},
\]
the heat comparison times obey
\[
  \tau_n
  =
  \frac{r_0^2}{\nu}4^{-n},
  \qquad
  \sum_{n=0}^{\infty}\tau_n
  =
  \frac{4r_0^2}{3\nu}
  <
  \infty.
\]
This finite arithmetic permits a schedule, but it does not establish one field realizing it.
If one \NavierStokes{} field exhibits localized structures at widths
\[
  r_0>r_1>r_2>\cdots\longrightarrow0
\]
and at sampled times
\[
  t_0<t_1<t_2<\cdots\uparrow T_*<\infty,
  \qquad
  t_{n+1}-t_n\asymp\tau_n,
\]
with comparison constants independent of \(n\), then that possible history has remaining time
\[
  T_*-t_n
  \asymp
  \sum_{k=n}^{\infty}\tau_k
  =
  \frac{4r_n^2}{3\nu}.
\]
At the \(n\)-th sampled width, the next gap and the total time left are both of order \(r_n^2/\nu\), and these gaps collapse as \(t\uparrow T_*\).
Even along this possible history, the first and successive Eulerian time derivatives of the same field at one fixed spatial point remain uncontrolled.
